\lecture{12}{Thu 23 Oct 2025 12:31}{Conformal diagrams}

We use the change of variables $u=t-r$, $v=t+r$, $U=\arctan u$, and $V=\arctan v$. Then
\[
	\mathrm{d} s^2=-\mathrm{d} u\mathrm{d} v+\left( \frac{v-u}{2}  \right) ^2\mathrm{d} \Omega =-\left( \frac{\mathrm{d} U}{ \cos ^2U} \right) \left( \frac{\mathrm{d} V}{\cos ^2V}  \right) +\frac{\left( \tan V-\tan U \right) ^2}{4} \mathrm{d} \Omega ^2
.\]
Trig identities give
\[
	=\frac{1}{\cos ^2U \cos ^2V} \left( -\mathrm{d} U\mathrm{d} V+\sin ^2(V-U)\mathrm{d} \Omega ^2 \right) 
.\]
We can draw cool diagrams.

\chapter{Einstein's equation}

We will now shift our focus. We will discuss Einstein's equation, which tells us what the metric is \emph{as a function of time}. This is much richer since the metric has to be governed by the existence of matter and energy, and the metric should respond to the existence of this.
